% ==========================================================
% EE 3150: Signals and Systems -- Lab 4 Report Template
% Fourier Series, Filtering, and Distortion
% ==========================================================
% Grading (52 pts): Intro 1 | Part 1: 15 | Part 2: 14 | Part 3: 15
%                   Conclusion 4 | Figure style 3
% Submit ONE typed PDF per team via Brightspace.
% Handwritten derivations may be scanned (see \includegraphics hints).
% ==========================================================
\documentclass[11pt]{article}

\usepackage[margin=1in]{geometry}
\usepackage{amsmath,amssymb}
\usepackage{graphicx}
\usepackage{float}
\usepackage{booktabs}
\usepackage{listings}
\usepackage{xcolor}
\usepackage{enumitem}
\usepackage{caption}
\usepackage{subcaption}
\usepackage[colorlinks=true,linkcolor=blue,urlcolor=blue]{hyperref}

% ---------- Code listing style ----------
\lstdefinestyle{pycode}{
  language=Python,
  basicstyle=\ttfamily\small,
  keywordstyle=\color{blue!70!black},
  commentstyle=\color{green!40!black},
  stringstyle=\color{red!60!black},
  numbers=left,
  numberstyle=\tiny\color{gray},
  frame=single,
  breaklines=true,
  showstringspaces=false
}
\lstset{style=pycode}

% ---------- Handy macros ----------
\newcommand{\todo}[1]{\textcolor{red}{[TODO: #1]}}   % delete before submitting
\newcommand{\Hjw}{H(j\omega)}
\newcommand{\Hjk}{H(jk\omega_0)}
\newcommand{\figplaceholder}[1]{%
  \fbox{\parbox{0.9\linewidth}{\centering\vspace{1.2cm}#1\vspace{1.2cm}}}}

\title{\textbf{EE 3150: Signals and Systems}\\[4pt] Lab 4: Fourier Series, Filtering, and Distortion}
\author{Colden Briggs \and Eben Radocchia \and Liam Donlan\\[4pt]}
\date{\today}

\begin{document}
\maketitle

% ==========================================================
\section{Introduction}  % [1 pt]
% ==========================================================
\todo{State the objective of the lab in a short paragraph: how Fourier-series
coefficients determine signal shape and sound (Part~1); what an LTI system does
to the coefficients (Part~2); what a nonlinear system does to them (Part~3).}

% ==========================================================
\section{Part 1 --- Building Instruments from Fourier Coefficients}  % [15 pts]
% ==========================================================

\subsection{Task 1.1 --- Musical Tones with Different Timbre}  % [7 pts]

\paragraph{Completed code} (3 pts) Include the completed \texttt{additive}
function and the four coefficient lists.
\begin{lstlisting}
import numpy as np
import matplotlib.pyplot as plt

fs = 8192
f0 = 220.0
t = np.arange(0, 2.0, 1.0/fs)

def additive(t, f0, amps, fs):
    x = np.zeros_like(t)
    for m in range(len(amps)):
        k = m + 1
        if k*f0 < fs/2:
            x = x + ...   # TODO
    return x

pure    = ...  # TODO
odd_1k  = ...  # TODO
all_1k  = ...  # TODO
all_1k2 = ...  # TODO
\end{lstlisting}

\paragraph{Stem plots of $A_k$ and brightness comments} (1 pt plot + 1 pt comment)
\begin{figure}[H]
  \centering
  \figplaceholder{Stem plots of $A_k$ vs.\ $k$ for the four tones\\
  \texttt{\textbackslash includegraphics[width=\textbackslash linewidth]\{fig\_stems.pdf\}}}
  \caption{Harmonic amplitudes $A_k$ for the four tones.}
  \label{fig:stems}
\end{figure}

\begin{itemize}[leftmargin=*]
  \item \textbf{Pure tone:} \todo{one-line comment on brightness}
  \item \textbf{Odd $1/k$:} \todo{comment}
  \item \textbf{All $1/k$:} \todo{comment}
  \item \textbf{All $1/k^2$:} \todo{comment}
\end{itemize}
\todo{Explain the trend: which harmonic content makes a tone bright, which makes it hollow or dull?}

\paragraph{One period of each tone} (1 pt plot + 1 pt pattern)
\begin{figure}[H]
  \centering
  \figplaceholder{One period of each tone, normalized to unit amplitude, same axes}
  \caption{One period of each tone (normalized to unit amplitude).}
  \label{fig:periods}
\end{figure}
\todo{Explain how the coefficient pattern defines the shape; which profile looks
like a square wave and which like a sawtooth?}

\subsection{Task 1.2 --- Envelope}  % [4 pts]

\paragraph{Completed code} (2 pts)
\begin{lstlisting}
def envelope(t, t_a=0.01, tau_d=0.15):
    return ...  # TODO
\end{lstlisting}

\paragraph{Click description and explanation} (2 pts)
\todo{Describe the audible difference between $t_a = 10$~ms and $t_a = 0$ for the
\texttt{all\_1k} tone, and explain why the abrupt onset produces a click
(think: discontinuity $\Rightarrow$ high-frequency content).}

\subsection{Task 1.3 --- Play a Melody}  % [4 pts]

\paragraph{Completed code} (3 pts)
\begin{lstlisting}
freqs = {"C4":261.63, "D4":293.66, "E4":329.63,
         "F4":349.23, "G4":392.00, "A4":440.00}
song = [("C4",0.5),("C4",0.5),("G4",0.5),("G4",0.5),("A4",0.5),("A4",0.5),
        ("G4",1.0),("F4",0.5),("F4",0.5),("E4",0.5),("E4",0.5),("D4",0.5),
        ("D4",0.5),("C4",1.0)]

my_amps = ...  # TODO
clip = []
for name, duration in song:
    tn = np.arange(0, duration, 1.0/fs)
    note = ...  # TODO
    clip.append(...)  # TODO
clip = np.concatenate(clip)
clip = clip/np.max(np.abs(clip))
# play clip at fs
\end{lstlisting}

\paragraph{Comment} (1 pt)
\todo{State which instrument ($A_k$) you used and explain in one line why the
melody is still recognizable.}

% ==========================================================
\section{Part 2 --- Filtering the Square Wave through the RC Network}  % [14 pts]
% ==========================================================

\subsection{Task 2.1 --- The Truncated Square Wave and the Gibbs Phenomenon}  % [4 pts]

\paragraph{Completed code} (1 pt)
\begin{lstlisting}
A = 1.0
f0 = 10.0
w0 = 2*np.pi*f0
T0 = 1.0/f0

def xtilde(t, N):
    x = np.zeros_like(t)
    for n in range(1, N+1):
        k = 2*n - 1
        x = x + ...  # TODO
    return x

t = np.arange(-0.25*T0, 1.25*T0, T0/10000)
# TODO: plot xtilde(t,N) for N in [1,3,7,49] and the ideal square wave
\end{lstlisting}

\paragraph{Partial-sum plot} (1 pt)
\begin{figure}[H]
  \centering
  \figplaceholder{$\tilde{x}_N(t)$ for $N=1,3,7,49$ with ideal square wave}
  \caption{Partial sums $\tilde{x}_N(t)$ of the square wave ($A=1$, $f_0=10$~Hz).}
  \label{fig:gibbs}
\end{figure}

\paragraph{Overshoot measurement and explanation} (2 pts)
\begin{center}
\begin{tabular}{@{}ccc@{}}
\toprule
$N$ & Peak value after jump & Overshoot above $A=1$ (\%) \\
\midrule
49  & \todo{} & \todo{} \\
199 & \todo{} & \todo{} \\
\bottomrule
\end{tabular}
\end{center}
\todo{Does the percentage shrink or stay roughly the same as $N$ grows? What does
this say about representing a discontinuous signal with finitely many harmonics?}

\subsection{Task 2.2 --- Filtering the Square Wave Using an RC Circuit}  % [10 pts]

\paragraph{Frequency response derivation} (1 pt)
Given $h(t) = \dfrac{G}{\tau_1-\tau_2}\left(e^{-t/\tau_1}-e^{-t/\tau_2}\right)u(t)$
with $\tau_1 = 10$~ms, $\tau_2 = 4.7$~ms, $G = 5.7$:
\begin{align}
\Hjw &= \int_{-\infty}^{\infty} h(\tau)\,e^{-j\omega\tau}\,d\tau \nonumber\\
     &= \todo{show steps} \nonumber\\
     &= \frac{G}{(1+j\omega\tau_1)(1+j\omega\tau_2)}.
\end{align}

\paragraph{Magnitude, phase, and comparison with Lab 3} (1 pt)
\begin{align}
|\Hjw| &= \todo{} , &
\angle \Hjw &= \todo{} .
\end{align}
\todo{Compare with Lab 3 and explain why they agree.}

\paragraph{Fourier-series output code} (2 pts)
\begin{lstlisting}
tau1, tau2, G = 10e-3, 4.7e-3, 5.7

def H(f):
    jw = 1j*2*np.pi*f
    return ...  # TODO

def y_fs(t, f0, N=99):
    w0 = 2*np.pi*f0; y = np.zeros_like(t)
    for k in range(1, 2*N, 2):
        Hk = H(k*f0)
        # TODO: compute the k-th harmonic of the output
        y = y + ...  # TODO
    return y
\end{lstlisting}

\paragraph{Convolution output code and overlay} (2 pts)
\begin{lstlisting}
def y_conv(f0, dt=5e-5):
    nper = 6
    t = np.arange(0, nper/f0, dt)
    x = np.sign(np.sin(2*np.pi*f0*t)); x[x == 0] = 1
    th = np.arange(0, 12*tau1, dt)
    h = ...   # TODO
    yc = ...  # TODO: np.convolve(...)*dt, keep first len(t) samples
    return t, yc
\end{lstlisting}

\begin{figure}[H]
  \centering
  \figplaceholder{Overlay of Fourier-series output and convolution output ($f_0=10$~Hz)}
  \caption{Fourier-series output vs.\ convolution output at $f_0 = 10$~Hz.
  The two agree once the initial transient has decayed.}
  \label{fig:overlay}
\end{figure}

\paragraph{Plots for $f_0 \in \{2, 10, 40\}$~Hz} (2 pts)
\begin{figure}[H]
  \centering
  \begin{subfigure}[b]{0.9\linewidth}
    \centering
    \figplaceholder{$f_0 = 2$~Hz: input and output, two periods}
    \caption{$f_0 = 2$~Hz}
  \end{subfigure}\\[6pt]
  \begin{subfigure}[b]{0.9\linewidth}
    \centering
    \figplaceholder{$f_0 = 10$~Hz}
    \caption{$f_0 = 10$~Hz}
  \end{subfigure}\\[6pt]
  \begin{subfigure}[b]{0.9\linewidth}
    \centering
    \figplaceholder{$f_0 = 40$~Hz}
    \caption{$f_0 = 40$~Hz}
  \end{subfigure}
  \caption{Square-wave input and filtered output over two periods.}
  \label{fig:threef0}
\end{figure}

\paragraph{Mathematical explanation} (2 pts)
\begin{center}
\begin{tabular}{@{}cccc@{}}
\toprule
$f_0$ (Hz) & $|H(j\omega_0)|$ & $|H(j3\omega_0)|$ & $|H(j5\omega_0)|$ \\
\midrule
2  & \todo{} & \todo{} & \todo{} \\
10 & \todo{} & \todo{} & \todo{} \\
40 & \todo{} & \todo{} & \todo{} \\
\bottomrule
\end{tabular}
\end{center}
\todo{For each $f_0$, explain which harmonics pass and which are attenuated, and how
that determines the output shape.}

\paragraph{Frequencies present in the output}
\todo{State which frequencies appear in the output relative to the input. This
statement will be tested against a nonlinear system in Part~3.}

% ==========================================================
\section{Part 3 --- Nonlinear Distortion}  % [15 pts]
% ==========================================================

\subsection{Task 3.1 --- The Clipped Waveform}  % [3 pts]
Parameters: $A=1$, $f_0 = 220$~Hz, $f_s = 44\,100$~Hz.

\begin{figure}[H]
  \centering
  \figplaceholder{One period of clipped sine for $L/A \in \{0.9, 0.5, 0.2\}$ with unclipped sine}
  \caption{Clipped sine $y(t)=\texttt{clip}(A\sin\omega_0 t,-L,L)$ for three clip levels.}
  \label{fig:clipped}
\end{figure}
\todo{Describe how the waveform changes as $L$ decreases. What shape does $y(t)/L$
approach as $L/A\to 0$?}

\subsection{Task 3.2 --- Analytical Fourier Coefficients of the Clipped Sine}  % [6 pts]
Let $\theta=\omega_0 t$ and $\theta_c=\arcsin(L/A)$. For odd $k$,
\begin{equation}
C_k = \frac{4}{\pi}\int_0^{\pi/2} y(\theta)\sin(k\theta)\,d\theta .
\end{equation}

\paragraph{Derivation (3 pts)} Split the integral at $\theta_c$.
\begin{align}
C_k &= \frac{4}{\pi}\left[\underbrace{\int_0^{\theta_c} A\sin\theta\,\sin k\theta\,d\theta}_{I_1}
      + \underbrace{\int_{\theta_c}^{\pi/2} L\sin k\theta\,d\theta}_{I_2}\right] \nonumber\\
\intertext{Using $\sin\theta\sin k\theta = \tfrac12[\cos((k-1)\theta)-\cos((k+1)\theta)]$
and $\cos(k\pi/2)=0$ for odd $k$:}
I_1 &= \todo{} \nonumber\\
I_2 &= \todo{} \nonumber
\end{align}
\todo{Or include a scanned handwritten derivation:
\texttt{\textbackslash includegraphics[width=\textbackslash linewidth]\{derivation.pdf\}}}

\paragraph{Final closed-form expressions (3 pts)}
\begin{align}
C_1 &= \todo{} \\
C_k &= \todo{}, \qquad k\ge 3 \text{ odd.}
\end{align}

\subsection{Task 3.3 --- Effect of Clip Severity}  % [3 pts]

\paragraph{Code (1 pt)}
\begin{lstlisting}
def clip_Ck(k, A, L):
    # handle k == 1 separately
    return ...  # TODO
\end{lstlisting}

\paragraph{Severity plot and description (1 pt)}
\begin{figure}[H]
  \centering
  \figplaceholder{$|C_k|$ vs.\ odd $k$ (up to 13) for $L/A \in \{0.9,0.5,0.2\}$}
  \caption{Harmonic magnitudes of the clipped sine for three clip levels.}
  \label{fig:Ck}
\end{figure}
\todo{Describe how the harmonic content changes as the clip deepens.}

\paragraph{Limiting cases (1 pt)}
\begin{enumerate}[label=(\arabic*)]
  \item $L/A\to 1$: $C_k \to \todo{}$; in particular $C_1\to \todo{}$.
  \item $L/A\to 0$: $C_k/L \to \todo{}$. This is the harmonic series of a
        \todo{waveform name}; it arises because \todo{explain from the shape in Task~3.1}.
\end{enumerate}

\subsection{Task 3.4 --- Listen, and Contrast with Filtering}  % [3 pts]

\paragraph{Audio (1 pt)}
\todo{Describe the clean vs.\ clipped tone synthesis ($f=220$~Hz, $A=1$, $L/A=0.5$,
envelope applied to both, $f_s=44\,100$~Hz, $\tau_d$ value used).}

\paragraph{What you hear}
\todo{Describe how the clipped tone sounds compared with the clean one and
connect it to the harmonics in Task~3.3.}

\paragraph{LTI vs.\ nonlinear (2 pts)}
\todo{In 2--3 sentences: state precisely what an LTI system does to the harmonics
of a periodic input, what a nonlinear system does, and why only one of the two can
produce frequencies absent from the input.}

% ==========================================================
\section{Conclusion}  % [4 pts]
% ==========================================================
\todo{Summarize the lab and highlight what you learned.}

\paragraph{Take-away messages}
\begin{enumerate}
  \item \todo{}
  \item \todo{}
  \item \todo{}
  \item \todo{}
\end{enumerate}

% ==========================================================
\appendix
\section{Complete Python Code (optional)}
% ==========================================================
% To include your script directly:
% \lstinputlisting{lab4.py}

\end{document}